\documentclass[../main.tex]{subfiles}
\begin{document}
%==========================================TIÊU ĐỀ TẬP========================================
\begin{table}[h]
  \begin{adjustwidth}{-1cm}{}
    \begin{tabular}{>{\centering\arraybackslash}p{6cm}|>{\raggedright\arraybackslash}p{14cm}}
      \multirow{5}{6cm}
      % Cột bên trái: ÉPISODE và số tập
      {\\[-40pt]
        \flaregothic\fontsize{20pt}{20pt}\selectfont \centering
        \color{eptype}HISTOIRE\color{black}\\
        \brian\fontsize{72pt}{72pt}\selectfont
        \color{epnum}6
        \\[18pt]
        % \flaregothic\fontsize{14pt}{14pt}\selectfont
        % \color{epattr}BIENVENUE, \uline{\color{Banpire} Mel} !
        \flaregothic\fontsize{18pt}{18pt}\selectfont
        \color{epattr} FIN DE TOME !
        \color{black}
      }
       &
      % Dòng thứ nhất: tên tập tiếng Nhật
      {\fotskip\fontsize{18pt}{18pt}\selectfont パーティーの\ruby{後}{あと}はお\ruby{泊}{と}まり\ruby{会}{かい}！} \\[6pt]
      % Dòng thứ hai: tên tập tiếng Anh
       & {\cafeteria\fontsize{18pt}{18pt}\selectfont After the Party Comes the Sleepover!}      \\[4pt]
      % Dòng thứ ba: tên tập tiếng Việt
       & {\vnmsans\fontsize{18pt}{18pt}\selectfont Bữa tiệc sau đại tiệc}                       \\[8pt]
      % Dòng thứ tư: Nhân vật xuất hiện
       & {\sen Track Used: Cometbound (\ruby{彗}{すい}\ruby{星}{せい}の\ruby{界}{かい}\ruby{隈}{わい})}
    \end{tabular}
  \end{adjustwidth}
\end{table}
%===========================================NỘI DUNG==========================================
\color{black}

Mười một giờ năm mươi lăm phút đêm.

Bên ngoài cửa sổ, thành phố Kawachi đã chìm dần vào giấc ngủ yên bình. Ngày sinh nhật của Kuon cũng đang khép lại, thời khắc đặc biệt ấy lặng lẽ trôi về cuối cùng của nó.

Công việc dọn dẹp lại tầng ba cuối cùng cũng hoàn tất. Căn phòng này, vốn được thiết kế như một không gian đa năng, giờ đây đã biến thành một khu trại tạm đầy sắc màu: bàn ghế được dọn gọn sang một bên, chăn gối và đệm trải kín sàn thạch cẩm, một vài chiếc túi ngủ màu mè do các thành viên mang theo từ nhà được đặt ở các góc phòng. Những dải đèn dây LED lung linh dọc theo trần nhà vẫn chưa tắt, tạo nên một bầu không khí vừa ấm cúng vừa rộn ràng.

Kuon cúi đầu xuống, nét mặt đầy chân thành: ``\dots Một lần nữa, anh cảm ơn mấy đứa đã đến dự sinh nhật anh nha! Để anh dẫn mọi người xuống cửa-''

Thế nhưng, khi cậu vừa định bước về phía cầu thang, một giọng nói quen thuộc vang lên, nhẹ nhàng mà đầy tinh nghịch: ``Ai bảo là bọn em sẽ về?''

Cậu quay đầu lại. Một nụ cười rạng rỡ đang hiện rõ trên khuôn mặt của Sora.

``Hả? Nhưng mà giờ đã muộn rồi mà? Gần nửa đêm rồi đấy\dots''

A-chan bước tới, tay cầm một bảng ghi chép như mọi khi, nhưng nụ cười của cô thì không còn nghiêm nghị như lúc làm việc thường ngày: ``Thực ra thì\dots\ mọi người quyết định sẽ ở lại đây ngủ đó, Hikawa-san.''

Cậu quản lý đứng sững tại chỗ, mắt mở to hết cỡ. ``\textbf{Ể!!?}'' Một tiếng thốt lớn vang lên không kiểm soát. ``\textbf{N-Nhưng\dots\ sao lại\dots?}''

Fubuki búng tay cái tách, chọc thêm ``Chà chà\~{} Bọn em chuẩn bị tinh thần cắm trại tại đây hết rồi nha, đừng nghĩ đến chuyện đuổi về!''

Pekora bật cười, tay ôm gối ngủ hình cà rốt: ``Anh nghĩ bọn em tới chơi không mang theo đồ ngủ sao, peko? Đánh giá thấp quá đó nha!''

``Với cả anh nghĩ bọn em về giờ này sẽ cực kỳ nguy hiểm sao?'' Rushia chỉ về phía ngoài, trời đã tối đen như mực.

``\dots Cũng đúng. Anh biết là Kawachi an toàn rồi, cơ mà anh nghĩ là mấy đứa nên-'' cậu định lên tiếng giải thích, nhưng lại bị Matsuri chặn họng tiếp.

Cô nàng nghiêng đầu, chỉ tay ra sau lưng cậu: ``Anh không thấy bọn em đã trải đệm sẵn rồi à?''

Kuon quay lại nhìn phía sau mình, đúng là sàn nhà giờ đã lấp đầy bởi một rừng chăn gối. Nhưng cậu vẫn lắc đầu, bối rối: ``Không\dots\ lúc trước anh không thấy gì cả\dots''

Lamy chớp mắt rồi giải thích với giọng vui tươi: ``Thì bởi lúc đó anh còn đang trên bếp dọn dẹp với mọi người mà! Bọn em đã tranh thủ làm mọi thứ trong lúc anh vắng mặt đó\~{}''

``Nhưng mà, tại sao chứ!?'' cậu vẫn chưa hoàn toàn hiểu nổi tình hình.

A-chan thở nhẹ một cái, chỉnh lại cặp kính, rồi kể lại đầu đuôi: ``Hồi nãy, lúc cậu đang say sưa hát - à, tức là sau khi bị bọn tôi ép lên sân khấu đó - thì Sora-chan đã nghĩ ra thêm một ý tưởng. Cô ấy nói rằng là\dots''

``\dots Nếu hư sinh nhật này là kỷ niệm đầu tiên của Kuon-senpai cùng với toàn bộ \hll, thì nhất định phải kết thúc theo cách đáng nhớ nhất đó!'' nàng Thần tượng hí hửng giải thích, trích gần như nguyên văn câu nói của mình ban nãy với cô bạn.

Tới đây, Kuon chỉ có thể thở dài, chẹp miệng một tiếng. ``Thế là cuối cùng mấy đứa định chơi tới bến luôn đấy à\dots'' cậu cười mếu, như đã định được số phận.

``Ban đầu tôi còn nghi ngờ, nhưng mà khi thấy Sora-chan cầu xin với đôi mắt long lanh\dots\ tôi đành thở dài và gật đầu đồng ý như cậu luôn.''

Cậu hướng về các cô gái, hít một hơi thật sâu. ``Như vậy là mấy đứa quyết không về hả?''

Mọi người đồng thanh, tinh nghịch và đầy hứng khởi: ``\textbf{Không! Chúng em sẽ ngủ lại tại nhà anh đó, Kuon-senpai!}''

Từ phía sau, bố mẹ Kuon cũng vừa bước lên tầng. Họ đã nghe thấy đoạn hội thoại náo nhiệt này từ xa.

Mẹ cậu cười rạng rỡ, còn không quên trêu thêm một câu: ``Tốt quá rồi còn gì? Được ngủ với hơn ba mươi người con gái thế kia thì đúng là thiên đường rồi con!''

``Thiên đường gì mẹ ơi\dots'' cậu thở dài, một tay chống hông, tay kia vò tóc như thể đang cố vắt óc tìm lối thoát khỏi tình cảnh trớ trêu. ``Thiên đường gì mà toàn là bẫy rập\dots\ Ai mà ngủ nổi chứ.''

``Như thể Onii-san có một dàn harem vậy!'' Ryuuhou chẳng bỏ qua cơ hội, giơ tay lên má giả bộ đỏ mặt và nhái giọng một nhân vật trong anime một cách đầy mỉa mai: ``Kyaa\~{} Senpai, đừng nhìn em bằng ánh mắt đó đi\~{}!''

``Ngoại trừ việc harem này không nằm trong kịch bản,'' Kuon lườm em gái, mặt đầy chán nản, ``và anh cũng không hề ký hợp đồng tham gia vở diễn này luôn.''

Không khí bỗng trở nên rộn ràng hơn hẳn khi tất cả cùng phá lên cười. Sự náo động ấy như xóa tan mỏi mệt của buổi tiệc, để lại một sự gần gũi khó tả, như thể tất cả mọi người đã thật sự gắn kết lại, trở thành một đại gia đình nhà Hikawa.

Cậu Tổng quay sang bố mẹ, định phản bác gì đó, nhưng rồi lại thôi. Cậu thở dài, nở một nụ cười khổ lẫn bất lực: ``Ờm, thôi thì\dots\ Con vẫn sẽ ngủ trên kia ạ. Mấy đứa cứ ngủ ngon-''

Nhưng chỉ vừa nhấc chân bước về phía cầu thang, một giọng nói ngọt ngào mà nghiêm nghị bất ngờ vang lên từ phía sau, khiến sống lưng cậu lạnh toát.

Từ Sora.

``Kuon-senpai\~{}'' cô nàng mỉm cười, nhưng ánh mắt thì trông không hề dễ chịu chút nào.  ``Khi em nói là \emph{bọn em sẽ ngủ ở đây}, thì ý là \emph{anh cũng sẽ ngủ ở đây với bọn em đó.}''

Cậu chưa kịp phản ứng, các cô gái trong căn phòng cũng lần lượt lên tiếng hùa theo:

``Đúng đó!!''

``Chịu trách nhiệm với bản thân đi!''

``Đừng có mơ trốn!''

``Đó không phải là lời mời đâu!''

``Vào đi anh ơi!''

``Anh mà trốn là mất điểm nặng lắm đó!''

``Kuon-senpai bị bắt rồi nhé\~{}''

Cứ thế, hàng chục giọng nói đồng loạt vang lên như một đợt sóng vỗ, đẩy cậu lùi lại từng bước và tiến tới gần hơn tới căn phòng tầng ba kia.

Cậu đặt tay lên trán, gục đầu xuống thành cầu thang như thể chịu đòn trời giáng, rồi quay sang cầu cứu: ``B-Bắt buộc thật á\dots?''

Ryujin gật đầu không chút do dự, tay khoanh lại như đang tuyên án: ``Đúng rồi! Con là chủ nhà mà, phải có tinh thần hiếu khách chứ!''

Ngay lập tức, cả ông bố, bà mẹ và cô em gái của Kuon phối hợp nhịp nhàng như một tổ đội kỳ cựu, đồng loạt đẩy cậu một cú thật mạnh vào giữa căn phòng tràn ngập chăn gối và tiếng cười của các cô gái.

``Cứ ngủ với bọn họ đi, con à!''

``Ơ!?''

Tới đây, cậu cảm giác như bị phản bội, nhưng cũng không thể làm gì khác ngoài việc ngoan ngoãn ngồi xuống giữa vòng vây của các cô gái. Không khí trong phòng lập tức trở nên náo nhiệt hơn bao giờ hết: tiếng cười nói, tiếng chăn gối xào xạc, và cả những lời trêu chọc không ngớt vang lên.

Kuon chỉ biết ôm đầu than thở: ``Mọi người\dots\ thật sự nghiêm túc đấy à\dots?''

Sau lưng cậu, cánh cửa đóng sầm lại. Mẹ cậu và em gái Ryuuhou đồng thanh một câu, giọng ngọt lịm như kẹo mà lại khiến sống lưng cậu lạnh toát: ``Chúc ngủ ngon nha, Kuon/Onii-san!''

*RẦM!*

Im lặng. Căn phòng tầng ba giờ đây hoàn toàn khép kín, và Kuon thì đứng trơ trọi giữa một biển người.

Cậu quay mặt lại, chỉ thấy một ``rừng'' con gái đang cười híp mắt, tay vẫy vẫy gọi cậu lại gần. Người thì đập tay lên đệm mời, người thì giăng chăn chừa sẵn chỗ trống, người thì\dots\ thậm chí còn ác liệt hơn, vẫy gối dọa đánh nếu cậu không lại mau.

``Thật luôn\dots?'' cậu thốt lên như đang trò chuyện với một thế lực vô hình nào đó, mong chờ một sự can thiệp từ trên cao. Nhưng không có gì đáp lại ngoài sự réo gọi nhiệt liệt của các thành viên \hll.

``Kuon-senpai! Chỗ này nè!''

``Anh ngủ chung với bọn này đi!''

``Đừng có ngó nghiêng gì nữa, ngồi xuống đây nào!''

``Không được nghiêng người về bên kia nha\~{}!''

``Đừng ngủ nghiến răng đó!''

Cậu nhìn từng gương mặt, toàn những người cậu quý mến, tôn trọng\dots\ và cả đôi phần sợ sệt nữa. Thứ sợ hãi từ họ không đến từ nguy hiểm, mà từ sự thật rằng dù có chống cự thế nào, cậu cũng chẳng bao giờ thắng nổi họ. Không khác gì châu chấu đòi đá xe vậy, khi một con châu chấu nhỏ bé không đơn thuần phải đá một xe, mà đến tận ba mươi cái, càng khiến cho cơ hội của nó đã dê-rô nay còn về hẳn đến mức âm.

Kuon nở một nụ cười chua chát. Cậu biết rõ rằng đây không phải là một trận chiến có thể dùng lý lẽ hay sức lực để thắng. Đây là một ``cuộc nội chiến'' mà phía bên kia có đến ba mươi người, toàn bộ tập thể \hll.

``\textit{Đêm nay sẽ là đêm mất ngủ đây\dots}'' cậu lẩm bẩm như lời cáo chung cho tự do của mình, rồi ngậm ngùi bước đến cái đệm giữa phòng như thể đang tiến ra pháp trường. Nhắm mắt, xuôi tay.

Không có cửa thắng. Không bao giờ.

\ct

Mặc dù đã xuôi tay đầu hàng số phận, Kuon vẫn hy vọng thầm rằng mình có thể lặng lẽ nằm xuống một góc, đắp chăn kín đầu và ngủ cho qua chuyện.

Nhưng không. Đó chỉ là ảo tưởng đẹp đẽ, còn thực tế thì phũ phàng hơn nhiều. Ngay khi cậu đặt đầu xuống gối, một vài tiếng thốt lên, có người thậm chí còn lao tới cậu.

``\textbf{Kuon-senpai! Dậy chơi trò truth or dare (thử thách hay sự thật) đi!}''

``Không chơi là tụi em lấy bút lông vẽ mặt đó nha\~{}''

``Hay là chơi trò kể chuyện ma đi! Ai thua phải chui vào chăn ngủ một mình ngoài ban công!''

Trước khi cậu kịp phản ứng, đã bị kéo bật dậy bởi ba người cùng lúc: Fubuki nắm tay phải, Nene nắm tay trái, còn Matsuri thì đẩy lưng, miệng thì hét: ``\textbf{Vẫn còn quá sớm để ngủ đó!}''

Kuon thở dài ngán ngẩm với độ nhây của các thành viên. ``Rồi, rồi, chơi thì chơi. Nhưng mà một tiếng sau phải lên giường ngủ đó.''

``\textbf{Vâng\~{}}'' tất cả mọi người cùng đồng thanh, hớn hở vì cuối cùng có được một bằng hữu.

``\dots Và nhỏ tiếng thôi. Mười hai giờ đêm rồi đấy.''

``V-Vâng\~{}''

\ct

Bàn tròn giữa phòng nhanh chóng hình thành. Đèn chính được tắt, chỉ còn lại vài đèn ngủ tỏa ánh sáng mờ mờ ảo ảo, như thể đang tạo không khí cho một buổi lễ triệu hồi thay vì một buổi tiệc ngủ. Mùi trà nóng, bánh quy, và\dots\ mùi rắc rối đang len lỏi vào không khí.

``Chúng ta chơi truth or dare đi!!, peko!'' Pekora hét lên như một vị thần chiến tranh vừa ban sắc lệnh.

``Ủng hộ!'' Fubuki giơ tay. ``Mà nếu ai không chơi thì\dots''

``\dots thì hôm sau sẽ thấy clip dìm hàng của mình lan truyền trong group nội bộ,'' Mio tiếp lời, mặt không đổi sắc.

Kuon khoanh tay ngồi ở mép bàn, ánh mắt lạnh tanh như muốn xuyên thủng không khí hỗn loạn. ``Anh chỉ tới đây để ngồi yên một góc và tự hỏi mình tại sao lại sống đến giờ này.''

``Kuon-senpai\~{}'' Matsuri đặt tay lên vai cậu với nụ cười ám muội, ``không tham gia là vi phạm nội quy tiệc ngủ \hll\ đó nha\~{}''

``Vâng, nội quy quái quỷ này vừa mới lập ra từ ba giây trước mà thậm chí còn chẳng có điều lệ cụ thể nào,'' cậu đảo mắt, giọng hàm ý mỉa mai. ``Còn nội quy nào nữa thì chắc phải lập ra luôn quy định về phòng cháy chữa cháy đi.''

Nhưng dù có than trời cỡ nào, cậu cũng không thoát được. Và thế là chiếc chai rỗng bắt đầu quay.

Trong lượt quay chai đầu tiên, chiếc nắp hướng về Sora.

``Chị chọn được nói thật,'' cô nàng mỉm cười.

``Vậy chị từng đem lòng thương mến ai trong công ty chưa?'' Polka hỏi, ánh mắt gian xảo không giấu được.

``Không trả lời cũng được mà ha?'' Sora nhìn quanh, tìm cách tránh câu hỏi.

``Không. Không được,'' Cả phòng đồng thanh như một dàn đồng ca đã luyện tập trước.

``Em là người đã chọn nói thật mà, đúng chứ?'' cậu Tổng chỉ ra. ``Nói thật và chính xác vào.''

Sora thở dài, gật đầu chấp nhận. ``\dots Thôi được. Có. Nhưng người đó bây giờ là bạn thân của mình nên mình giấu luôn từ lâu rồi.''

Toàn phòng ré lên, còn Kuon thì nhìn cảnh tượng đó, lẩm bẩm: ``\textit{Chỉ hy vọng em ấy không chỉ đến mình\dots\ Mà hình như em ấy đang nói đến A-chan nhỉ?}''

Tới lượt quay thứ hai. Chiếc đầu chai chỉ thẳng về phía Kuon.

``Nói thật hay dám làm đây, Kuon-senpai?'' Subaru hỏi, miệng nhai bánh quy, tay vẫn lia điện thoại quay phim.

``Thật,'' cậu đáp gọn.

``Anh đã từng mơ thấy ai trong số tụi em chưa, peko?'' Pekora nheo mắt hỏi, giọng như thể chuẩn bị đâm một phát chí mạng.

Cả phòng vỡ oà với tiếng hú hét cổ vũ, ``Trả lời điiii!'', ``Chọn tên ra luôn!''

Kuon không do dự, đáp thẳng: ``Có.''

``\textbf{Ểểểểể!?}'' câu trả lời quá thẳng thắn khiến một số người suýt sặc trà.

``K-Không ngờ là anh ấy lại không hề tỏ vẻ giấu nhẹm đi đấy\dots'' Towa sững sờ.

``Dù gì đó cũng là tính cách của cậu ta mà,'' A-chan đáp.

``Cơ mà tại xao lại nư thế, Kuonsa-senpai'' Luna long lanh đôi mắt hỏi. ``Có phải vì bọn em dễ thưn quá hơm?''

``Không phải.''

Lập tức, một bầu không khí im lặng tỏa ra khắp căn phòng, như thể họ vừa phải nghe một thông tin chấn động.

``V-Vậy là tại sao?'' Fubuki hỏi, không giấu nổi vẻ ngạc nhiên.

``Phần lớn là do mấy đứa nghịch quá, làm nổ văn phòng, dọa khách mời, làm rơi đạo cụ, đốt giấy tờ\dots\ nên hình ảnh mấy đứa cứ ám vào giấc mơ của anh như hồn ma. Cứ tưởng bị stress công việc, hóa ra là mấy đứa tạo nghiệp chứ không phải anh tự tưởng tượng.''

Lại một sự im lặng khác.

Matsuri há miệng, Fubuki cố nhịn cười, Mio lẩm bẩm: ``Biết ngay kiểu gì anh cũng sẽ nói vậy\dots'' còn Pekora ngồi chết lặng.

Nene thì\dots\ đã lăn ra cười, mặc dù câu nói của cậu cũng chỉ thẳng cả chính cô nữa.

``Tôi hiểu cảm giác này của cậu,'' nàng Đeo kính tiến tới, vỗ vào vai cậu an ủi. ``Gần như đêm nào tôi cũng phát rồ vì giấc mơ điên loạn từ mấy đứa đó đấy.''

Kuon chống tay lên bàn, giọng bình tĩnh đến đáng sợ: ``Thực ra còn có một đêm anh còn mơ thấy bị khóa trái trong văn phòng cùng với ba người, Pol-pol, Marine-chan và Noel-chan. Cả ba đều bảo anh phải\dots\ làm sao nhỉ, `giải quyết mâu thuẫn nội bộ.' Tỉnh dậy mà tim còn đập thình thịch, tưởng là tai biến luôn à!''

Marine vỗ bàn cười: ``Tuyệt đối là mơ đẹp đó nha!''

``Không. Mơ đó khiến anh phải đi uống trà hoa cúc ba ngày liền để ổn định tâm lý. Anh vẫn còn nhớ như in lúc em sán tới đòi cưa cẩm anh đấy. Môi của em hôi rình kinh dị luôn!''

Lần này thì cả phòng lăn ra cười, còn nàng Thuyền trưởng thì hờn dỗi vì bị Kuon nói xấu.

Lượt quay chai tiếp theo. Vẫn lại là cậu Tổng xui xẻo đó. ``Lại là anh à?''

Nhưng khác với ban nãy, cậu lại được\dots\ đặc cách được một người trong nhóm chọn sẵn: ``Chọn dám làm nhé, Kuon-senpai!'' Không ai khác là Fubuki, cô nàng Cáo trắng lắm mưu mẹo.

``Dám làm thì dám làm, sợ đếch gì,'' cậu nhún vai, có vẻ là chịu chơi, nhưng ánh mắt lại đang toan tính một âm mưu khác.

Pekora ranh mãnh không tiếc lời: ``Vậy thì giả giọng moe và nói `Mấy anh ơi đừng chọc em nữa, em nhột đó\~{}' đi, peko!''

Mọi người ré lên, Matsuri còn bật nhạc nền đáng yêu trên điện thoại. Cùng với đó là những tiếng hò reo thúc giục cậu làm thật nhanh.

``Làm đi, làm đi nào senpai!''

``Bọn em sẽ cố gắng bắt trọn khoảnh khắc này!''

``Kuon-chan đáng iu quá đi, ne!''

Kuon nhìn mọi người một vòng, hít sâu một cái, rồi\dots\ bắt tay vào công việc của mình. Cùng với một chút kế hoạch.

Với một giọng nam trầm chuẩn chỉnh, cậu nghiêng đầu và bắn ra câu nói đó, y chang yêu cầu từ họ\dots

\dots ngoại trừ việc nó không hề dễ mến chút nào, mà nghe như một lời nguyền đến từ một con ma nữ trồi lên khỏi một cái giếng.

``\dots Mấy anh ơi đừng chọc em nữa, em nhột đó\~{}\dots''

Để cho chắc kèo hơn, cậu tự dùng đèn pin điện thoại của mình rọi thẳng vào mặt, nhăn mặt lại cố tỏ vẻ đáng sợ.

Không khí ngưng đọng. Một giây, hai giây\dots

Rồi những tiếng la hét của mọi người vang lên, như thể họ nhìn thấy ma.

Chỉ riêng Polka bật cười lăn lộn: ``\textbf{Moe cái gì chứ!? Anh nói như thể đang ám lời nguyền lên bọn này ấy!}''

Subaru đập gối xuống sàn, run rẩy toàn thân: ``Đáng sợ quá, đáng sợ quá! Lạnh cả sống lưng rồi\dots''

Fubuki thì giơ tay xin thua: ``Được rồi\dots\ anh thắng\dots\ lần này thôi đó!''

Thấy rằng mọi người không còn ai lên tiếng thách thức nữa, Kuon nhếch môi, giọng điềm nhiên: ``Trả thù hoàn tất.''

Cậu chỉnh lại cổ áo áo ngủ như thể vừa kết thúc một cuộc họp cổ đông: ``Giờ anh đi ngủ được chưa?''

``Chưa xong đâu nha\~{}'' Rushia nói với nụ cười tươi như chưa hề vừa bị dọa.

``Trời ơi, lại còn gì nữa?''

``Kể chuyện ma!''
một loạt giọng nữ đồng thanh vang lên như dàn đồng ca ác mộng.

Kuon thở ra một hơi dài. ``Rồi rồi. Miễn đừng bảo anh mặc đồ trắng đội tóc giả để kể chuyện là được.''

``Tất nhiên là không rồi!'' Noel đáp. ``Chỉ là kể chuyện kinh dị bình thường thôi!''

\ct

Kuon quay lại chiếc bàn tròn, tham gia trò kể chuyện ma, một ý tưởng mà cậu nghĩ là cuối cùng cũng có gì đó nghiêm túc hơn chút. Đèn được hạ thấp thêm một nấc, và các thành viên tham gia hoạt động nhỏ đó - Pekora, Mio, Sora, Kuon, Matsuri, Lamy, Ayame, Kanata, Korone, Miko, Towa - chui rúc vào mền, tạo thành một hình tròn như thể đang thực hiện nghi lễ gọi hồn.

``Vậy chị bắt đầu kể chuyện của mình nhé,'' Sora nở nụ cười hiền hậu, chỉnh lại độ sáng của đèn sáng hơn đôi chút, rồi cầm lấy chiếc đèn ngủ đó, giả vờ làm cây nến.

Với gương mặt điềm nhiên, ánh mắt không chút cảm xúc và giọng nói trong veo, cô bắt đầu kể một câu chuyện về một chiếc tủ khóa trong phim trường luôn phát ra tiếng lạch cạch vào nửa đêm.

``Chuyện kể rằng trong studio chúng ta có một chiếc tủ khóa, đây cũng là nơi mà trang phục hay để ở đó.''

Tất cả mọi người tiếp tục im lặng lắng nghe.

``Một đêm nọ, khi chị là người ở lại cuối cùng theo yêu cầu của A-chan để kiểm tra lại đồ đạc - hồi đó Kuon-senpai chưa gia nhập hội chúng ta đâu - bỗng chốc\dots\ có những tiếng lạch cạch phát ra từ chiếc tủ khóa đó.''

Một vài thành viên nhát cáy, như Mio và Nene, bắt đầu run nhẹ ngay khi chi tiết tủ khóa được nhắc đến.

``Chị đã đi tới kiểm tra, nhưng bên trong đó\dots\ không có lấy một cái gì cả, ngoài bộ quần áo của chúng ta,'' Sora tiếp tục kể, giọng có phần căng thẳng hơn.

``Chác là bọn chuột lục lọi tìm đồ ăn thôi mà,'' Botan khoanh tay đáp, giọng tỉnh bơ.

``Ừ, dễ thế lắm. Anh thấy cũng bình thường mà,'' Kuon tiếp lời.

``Nhưng mà chưa hết đâu!'' nàng Thần tượng chợt cao giọng, khiến những người quanh chiếc bàn tròn đó khẽ giật mình. ``Sáng hôm sau, có thêm một bản vẽ storyboard (kịch bản phân cảnh) mới đặt trong đó. Không ai thừa nhận đã để nó vào cả.''

Tới đây, một vài người đã bắt đầu rùng mình trước sự xuất hiện kỳ quặc của tấm kịch bản đó. Nhưng về phần đông, trông họ vẫn còn khá tỉnh bơ, không có dấu hiện sợ sệt gì.

``Ể\dots\ Sora-senpai lại cái giọng đó nữa rồi!'' Mio rúc vào vai Okayu. ``Sao mặt chị ấy không thay đổi gì hết vậy trời!?''

``Đó mới là điều đáng sợ,'' Botan gật gù. ``Công nhận Sora-senpai cứng cỏi thật.''

``À không, chị thấy thế này cũng bình thường mà,'' Sora tươi tỉnh nói, khuôn mặt giãn ra, đặt lại chiếc đèn ngủ xuống bàn. ``Chị thấy có nhiều ca như thế ở công ty chúng ta rồi.''

Một số người nghe vậy liền nhíu mày, bán tín ban nghi với lời cô nói. Họ quay sang Kuon, tìm kiếm câu trả lời từ cậu Tổng.

``Đáng sợ cái gì chứ, em ấy kể đúng mà, chuyện đó xảy ra như cơm bữa ở huyện,'' Kuon chống cằm thở dài. ``Rõ ràng là ai đó đi làm khuya rồi bỏ quên storyboard thôi.''

``Lý do gì kỳ vậy, ne!?'' Miko nhăn mặt.

``Em tưởng là có mỗi anh với A-san tăng ca đêm chắc? Nhiều người cấp dưới bọn anh còn `chạy sô' đêm đến sấp mặt kìa,'' Kuon giải thích. ``Nên anh cho đó cũng bình thường thôi. Mấy đứa khác còn câu chuyện nào khác thú vị hơn không?''

Sau Sora, vài người khác cũng bắt đầu kể câu chuyện của mình. Có người thì kể câu chuyện do chính họ thêu dệt, người thì dựa theo một truyền thuyết đô thị nào đó mà có lẽ mọi người đều đã biết (hoặc chưa).

NMhưng đến lượt Mio thì mọi thứ bắt đầu\dots\ náo loạn. Trái với sự điềm tĩnh như mọi người, hoặc cùng lắm là hơi run sợ như Miko hay Lamy, nàng Sói đen này\dots\ đã run cầm cập ngay khi cô còn chưa bắt đầu câu chuyện.

Trong cơn run rẩy, cô nàng lắp bắp kể lại một câu chuyện kinh dị về một chủ đề khá muôn thuở: nhà vệ sinh.

``E-Em nghe nói\dots\ có một c-câu chuyện đáng sợ đã x-xảy ra\dots\ ở nhà vệ sinh tầng ba trong công ty chúng ta!!''

``Ý em là cái nhà vệ sinh dành cho bên nhân viên á? Vì anh không nhớ là có nhà vệ sinh nào cho thần tượng cả, nhờ vào lời giải thích của thiên tài nào đó\dots'' cậu thản nhiên nói, ngước nhìn về phía Miko, người đề ra ý tưởng không lắp toa-lét cho các thành viên, vì cô cho rằng ``thần tượng không bao giờ đi vệ sinh.''

``M-Miko tưởng là idol không cần đi vệ sinh chứ, ne!'' nàng Vu nữ giật thót, vội minh oan.

``Và tôi cũng giải thích rằng idol cũng là vật thể sống, Sakura-san,'' A-chan chỉ ra. ``Nhưng rồi cô vẫn nằng nặc không muốn xây chúng. Tôi vẫn đang cần câu trả lời vì sao đấy.''

``Ơ, là thì, vì\dots'' Miko luống cuống, tìm lời giải thích hợp lý cho yêu cầu có phần ngớ ngẩn và vô lý đó.

``D-Dù sao thì!'' Mio rít lên, cố gân lên để đưa mọi người về câu chuyện. ``C-Có một\dots\ một lần\dots\ mình nghe tiếng gọi\dots\ từ phòng cuối\dots\ mà không có ai bên trong hết!!!''

``Chắc là em mệt quá nên nghe nhầm thôi,'' Kuon lẩm bẩm, khẽ bật cười. ``Cụ thể là tiếng đó thế nào vậy?''

``Em còn nghe tiếng *khụ khụ*\dots\ như n-người già bị câm nữa chứ\dots!'' nàng Sói đen xanh mặt. ``Em còn nghe thấy tiếng lợn kêu khi đang bị chọc tiết trong đó nữa!''

``À, em đang nói đến Hitori-san chứ gì. Hôm trước lúc đang dọn dẹp sân khấu anh thấy chú ấy ngủ quên trong đó, ngáy ro ro như lợn kêu đó,'' giọng cậu bằng bằng, như thể đó là một câu chuyện bình thường khác.

``H-hả!? Thật hả!? Thế mà em cứ tưởng\dots''

``Anh còn lạ gì chú ấy nữa. Ông ấy hay trốn sếp bằng cách giả vờ ngủ trong toilet. Thế mà hôm đó ngủ thật luôn!''

Toàn bộ căn phòng ồ lên. Những tiếng cười rộ lên trước sự việc có phần oái oăm đó, còn Mio thì thở phào nhẹ nhõm vì cuối cùng bí ẩn tiếng động trong nhà vệ sinh kia cũng có lời giải đáp.

``Tôi sẽ báo cáo việc này lên YAGOO-san\dots'' A-chan lẩm bẩm, ghi lại lời chứng cứ gián tiếp từ cậu Tổng.

Đến lượt Kuon kể câu chuyện của mình. Ai nấy đều chờ mong hoặc là một câu chuyện đáng sợ thật sự để đẩy không khí cao trào lên, hoặc là một pha tấu hài khác để làm dịu không khí đi.

``Câu chuyện này anh cũng chỉ mới luận ra gần đây thôi,'' cậu bắt đầu dẫn dắt. ``Nó có thể không chính xác, nhưng đại ý thế này\dots''

Tất cả im phăng phắc. Có người đã vội nuốt nước bọt ngay khi nhìn thấy bản mặt có phần quá nghiêm túc từ cậu.

``Mấy đứa biết trong văn phòng chúng ta có một phòng thu âm, đúng chứ? Đó cũng là nơi mà đôi lúc anh ngủ lại để kiểm tra kịch bản cho các buổi podcast của mấy đứa đấy.''

Subaru và Korone gật đầu. Chính họ đã từng vào trong căn phòng đó một lần, với một người làm một buổi ASMR thảm họa\footnote{Xem lại {\flaregothic SPÉCIAL VI}, quyển số Bốn}, và một người đã từng lên ý tưởng phá căn phòng đó khi cô cố mở căn phòng bằng cửa sập\footnote{Xem lại {\flaregothic ÉPISODE 78}, quyển số Sáu.}

``Anh có nghe đồng nghiệp kể rằng trong đêm tối, nếu có quá nhiều người tụ tập trong căn phòng đó, một điều kỳ quặc sẽ xảy ra.''

Một vài người bắt đầu rúc sát vào nhau hơn. Không phải vì nội dung của nó quá dị (thực tế là có thể, nhưng Kuon chưa kể đến đoạn đó), mà bởi tông giọng trầm cùa cậu khiến cho họ bất an.

``C-Cụ thể là thế nào vậy, peko?'' Pekora khẽ cụp tai xuống, nhỏ giọng hỏi.

``Hừm\dots'' cậu khoanh tay lại, làm bộ lục soát lại trí nhớ, rồi đột ngột vỗ tay mạnh một tiếng làm cho họ ớn lạnh hơn. ``Anh nghe nói là hệ thống thu âm sẽ tự động khởi động, và ghi lại mọi tiếng động trong phòng. Từ những gì chúng ta đang nói ở đây, cho tới những thứ âm thanh khe khẽ, như tiếng thở, tiếng xê dịch đồ đạc, thậm chí là\dots\ tiếng mút kẹo nữa\dots''

``Ể!?'' Luna la lên, vội bỏ chiếc kẹo mút khỏi miệng. ``K-Khíp thế-nora!''

``Nhưng mà chưa hết đâu. Cái phần kinh dị hơn mới là thứ làm người ta khiếp đảm đây: đoạn ghi âm đó sẽ tự động gửi đến email của quản lý cấp cao - thường sẽ là anh, A-san hoặc Tanigo-san, nhưng đôi khi có thể là một người khác.''

``A-Ai co?'' Matsuri cố sát lại tới gần Aki, người vẫn đang chăm chú nghe câu chuyện.

``Sáng hôm sau, người nhận được email đó sẽ in nguyên văn những gì chúng ta nói, những tiếng động ta phát ra và dán chúng lên bảng thông báo nội bộ.''

``T-Thế cơ à?'' Subaru xanh mặt. ``N-Nhưng mà chúng ta có thấy được gì từ bản thông báo đâu?''

``Đó là bởi vì anh, A-san hoặc Tanigo-san đã kịp tiêu hủy chúng trước khi mấy đứa biết rồi. Nhưng anh vẫn còn giữ hai bản ghi lời đấy.''

``N-Nó như thế nào vậy?'' Lamy ôm chầm lấy Botan, ngước mắt đang rưng rưng lệ về cậu.

``Xem nào\dots'' cậu lấy ra tập tài liệu được chuẩn bị sẵn (thực chất là giấy trắng không viết gì) và cầm lấy cây đèn ngủ đọc lên: ``Một bản ghi chép từ đêm 21/3/2020, khi mà Subaru-chan ở trong căn phòng đó.''

``Ể-Ể!? E-Em đấy ạ!?'' nàng Vịt rít lên.

``Nguyên văn của nó đây: *tiếng xóc nước*, *tiếng xóc nước*, Subaru: A-S-M-R\dots\ Các bạn đang nghe thấy tiếng sóng vỗ\dots\ Tiếng sóng vỗ nhè nhẹ đang chảy mật vào tai các bạn đây\dots''

Ngay khi cậu đọc những dòng đó lên, nàng Tomboy vội lao tới cái chăn mà Okayu đã dựng lên từ trước và trốn dưới đó, chỉ hé đôi mắt xanh của mình lên.

``Anh còn nghe nói rằng nếu căn phòng đó có đúng tám người tụ lại ở một điểm\dots'' cậu bất chợt hạ thấp giọng, rồi ngước mắt hình viên đạn lên về phía mọi người, ``\dots thì hệ thống sẽ nhầm lẫn, tưởng là đang tổ chức buổi thử giọng bất hợp pháp. Và thế là\dots\ sẽ có người bị gạch tên khỏi lịch phát sóng trong hai tuần.''

Im lặng.

``\textbf{K-Không thể nào!!!}'' tất cả mọi người trong chiếc bàn tròn đó thất thanh.

``C-Chả trách mà mình toàn bị ăn ban (bị cấm túc) không à!'' Marine hoảng hốt. ``Chính em với cả hội và thế hệ thứ Tư cùng vào trong đó ăn vụng đó!''

``Tự dưng lại kéo thêm cả em vào làm gì chứ, Marine-senpai!?'' Kanata cũng không giấu nổi sự hốt hoảng. ``Giờ thì em sắp bị cấm túc rồi này\dots''

``Cháy nhà mới ra mặt chuột rồi à\dots'' A-chan thở dài, lặng lẽ khi tiếp tội vào trong sổ tay của mình.

Nhìn cảnh các thành viên của thế hệ thứ Ba và thứ Tư bát nháo khi chính họ vô tình vướng phải một tội khá vặt vãnh, để rồi bị câu chuyện kinh dị của Kuon thuyết phục, cậu chỉ có thể cười trừ.

``Thế ai ngồi gần Kuon nhất thì an toàn phải không!?'' Nene thốt lên, chỉ về phía cậu.

``C-Có thể lắm!''

Và như một phản xạ sinh tồn, bảy, tám người cùng lúc nhào vào cậu Tổng. Hết Lamy, Ayame, Kanata, rồi lại tới Korone, Miko, Towa, Matsuri\dots\ ai ai cũng chen chúc thành một đống hỗn loạn bên cạnh, thậm chí có người còn chồng hẳn lên cậu.

``T-Thôi mấy đứa! Anh không phải bệ thờ đâu!'' Kuon kêu lên. ``K-Khó thở quá!''

``V-Vậy có lời nguyền nào không?'' Lamy ép vai cậu, gặng hỏi.

Tới đây, cậu Tổng cười một tràng lớn, thấy rằng họ đã tin sái cổ vào câu chuyện của cậu: ``Câu chuyện anh bịa ra mà mấy đứa tin đến thế sao? Làm quái gì có chuyện đó!''

``C-Cái- Kuon-senpai đúng là!''

Cuối cùng, không những cậu được ``tặng'' những lời mắng nhiếc từ các cô gái, mà còn phải nhận thêm mấy cú đánh yêu từ họ vì cậu kể giống như thật.

``Anh nói thế có khác gì dọa người ta khỏi bữa tiệc đâu\dots'' Okayu vẫn còn đang cười lăn lộn.

``Thì đúng là thế còn gì nữa. Mấy đứa định đè chết anh ngay trong đêm hậu sinh nhật để lên bàn thờ ngắm gà khỏa thân chắc?''

Một vài thành viên trong nhóm khẽ bật cười với lời bình nửa đùa nửa thật kia.

\ct

Sau màn kể chuyện ma, một căn phòng đầy những chiếc gối, những tấm đệm futon bừa bộn cùng với ánh đèn mờ ảo của những chiếc đèn ngủ khiến không khí càng thêm mộng mị.

Một giờ sáng ngày Mười Tám. Ngày sinh nhật của cậu đã dần xa, nhưng cuộc vui vẫn còn tiếp tục.

Ngoại trừ nhân vật chính của bữa tiệc, cũng là người con trai duy nhất trong că phòng toàn là con gái và phụ nữ này.

Kuon lặng lẽ lăn qua một góc, trùm chiếc chăn lên tận mặt để không bị những thứ xung quanh phá đám cậu.

``\vie{Mệt quá, mệt quá\dots{} Ngủ thôi, mai còn đi làm\dots}'' cậu lẩm bẩm, hy vọng rằng cuối cùng cậu cũng có thể yên tâm đi ngủ. Cậu ráng nhắm mắt, cầu cho cơn buồn ngủ mau kéo tới giữa một bầy con gái vẫn còn rôm rả trò chuyện và làm những trò ngớ ngẩn.

Nhưng hỡi ôi, đời không hề như là mơ.

Ngay khi cậu vừa nhắm mắt và bắt đầu chìm dần vào giấc mộng\dots

*BỤP!*

\dots một chiếc gối bay thẳng vào mặt cậu, khiến cậu rên lên mạnh và bật dậy như lính trực chiến.

``Dậy đi nào, Kuon-senpai!'' một giọng nói reo lên đến từ Matsuri, người đang cầm một chiếc gối khác. ``Giờ thì không ai bị loại khỏi cuộc chiến gối nữa nhé!''

``Chiến gối\dots? Mấy đứa là trẻ con à?'' cậu lồm cồm bò dậy, quẳng cái gối vừa rồi ra xa, vô tình trúng phải Fubuki.

``Ái da\~{} Kuon-senpai cũng muốn chơi ném gối à?'' nàng Cáo trắng nhoẻn miệng cười.

``Không. Anh muốn đi ngủ. Mà mấy đứa cũng đi ngủ đi, một giờ sáng rồi đấy\dots''

``Hơm đực\~{} Zẫn còn quá sớm để ngủ mà-nora\~{}'' Luna cầm một chiếc gối nhỏ khác, vui vẻ chêm vào. ``Kuonsa-senpai cũng phải tham gia chứ?''

``Thôi anh xin, tha cho anh đi. Anh có phải cú đêm đâu\dots''

Ngay khi câu nói của cậu vừa dứt, một loạt gối bắt đầu bay như pháo bông khắp phòng. Không rõ ai là kẻ phát động, chỉ biết sau vài cú bốp bốp, mọi người đã hò hét ầm trời như đang mở hội.

``\textbf{Ném gối thôi!!!}''

Pekora ném trúng mặt Subaru, khiến cô nàng gào lên như vịt bị bóp cổ: ``Ặc!''

Fubuki núp sau đống đệm dựng lên như boong-ke, chĩa một cú chuẩn xác thần sầu thẳng vào gáy Mio đang lơ đễnh.

``Đoàng\~{}☆''

Lamy thì cười không ngớt, ôm gối như bom tạ, lừ lừ tiến lại như một con trùm mini giữa bãi chiến trường.

Nene trượt chân trên đệm và vô tình bay nguyên người vào Polka, hai người xoắn thành một cục rồi cùng rơi khỏi tấm nệm.

``Bùm!''

``N-Nhột quá, Nene! Bỏ tôi ra! Ahaha!''

Choco và Mel thì\dots\ không ném gối vào nhau, mà thay phiên ``bắn tim'' vào đối phương, khiến Mio giật thót tim vì tưởng trúng đạn thật.

``Nhận lấy này\~{}'' nàng Dơi ném chiếc gối bông về phía nàng Sói đen, làm cho cô bị ``hạ gục'' và lăn xuống đất. ``Đừng có đùa kiểu thế chứ!''

Okayu, người lại đang nằm trong giấc nồng, vừa cười vừa nói mớ trong mơ: ``Gối này có vị cơm nắm đây\~{} Ha ha\dots''

A-chan, một nhân viên kỳ cựu, đại diện cho lý trí và trật tự, lại\dots\ đang nằm ngay sát tường, ngủ ngon lành từ lúc kể chuyện ma xong, gối ôm ôm chặt, miệng còn lẩm bẩm ``\textit{deadline\dots\ meeting\dots\ không được trễ\dots}'' mà hoàn toàn không hay biết gì đang xảy ra xung quanh.

Kuon ở một chiến tuyến hoàn toàn khác, cậu là người duy nhất không muốn dây dưa vào cuộc hỗn chiến giữa những chiếc gối, mà chỉ cầu mong được ngủ để sáng mai còn đi làm. Cậu lặng lẽ rút lui chiến thuật trong lúc mọi người vẫn còn đang hăng say trong cuộc chiến, trườn khỏi hiện trường, chui vào một góc đệm gần cửa sổ, hít thở sâu.

Cậu nhắm mắt lại, đầu không nghĩ tới những gì đang xảy ra nữa.

Miệng cậu lẩm bẩm bắt đầu bài đếm cừu quen thuộc.

``\textit{Một con\dots\ Hai con\dots}''

Mọi thứ bắt đầu yên ắng lại. Cậu đang dần chìm vào giấc mộng một lần nữa\dots

*BỤP!*

Một chiếc gối khác, lần này có vẻ nặng hơn, đập thẳng vào lưng cậu.

Cậu choàng dậy, hét lớn: ``\textbf{\vie{Cái quái gì vậy!? Để yên cho anh ngủ cái!}}''

``Không đâu\~{} Đêm vẫn còn dài mà\~{}'' Nene ngân nga, ném thêm một chiếc gối khác vào đối thủ.

Nghe đến tông giọng nghịch ngợm của nàng ``Hải cẩu'' đang vui vẻ ném gối, cậu bắt đầu cảm thấy nóng mắt, tức đến mức muốn xì khói vì phá đám giấc ngủ.

Tuy thế, cậu vẫn cố gắng kiềm chế, hít một hơi thật sâu, nhẹ nhàng đặt đầu xuống chiếc gối một lần nữa. Để cho cẩn thận hơn, cậu còn đáp chăn kín mít cả mặt, hoặc là vì điều hòa lạnh, hoặc bởi cậu không muốn dính líu vào trò ném gối vô bổ đó nữa.

Không khí ấm dần. Cơn buồn ngủ ùa đến một lần nữa\dots

*BỐP!*

\dots chỉ để rồi lại bị ngắt quãng khi một chiếc gối khác bay trúng vào thái dương của cậu.

``Grừ\dots''

``\textit{Đứa nào!? Ai!? Vì sao!? Có công lý nào cho người lao động vục mặt vào guồng quay công viejc không!?}'' cậu cố nhẫn nhịn, hét lên trong lòng những điều đó.

Chỉ để rồi ngay khi cậu vừa ngẩng mặt dậy và định thốt ra những lời đó, một chiếc gối khác dày cộp như gạch bay thẳng vào mặt.

``Hây cha! Lại trúng Kuon-senpai rồi kìa!'' Noel hô lớn.

``Hay là chúng ta cứ tiếp tục tấn công anh ấy đi!'' Towa nhanh nhảu bày trò.

Tới đây, cậu cảm thấy giới hạn của mình đã đạt đỉnh.

Không còn gì để mất, cũng không còn lý do gì để kiềm chế cho những kẻ dám phá bĩnh giấc ngủ của cậu.

Và rồi, trong những tiếng cười nói của các cô gái dưới lúc đã quá giờ Sửu\footnote{Tính từ 1 giờ đến 3 giờ sáng}, cậu dồn hết sức lực của mình\dots

``\vie{\textbf{\Large{Chịu hết nổi rồi đó!!!}}}''

\dots một tiếng hét vang lên như sấm đánh, làm choàng tỉnh cả những con người vốn dĩ không liên quan gì tới trận đấu gối này (trừ A-chan, cô nàng vẫn đang ngủ ngon lành.)

Những người còn thức thì đông cứng tại chỗ. Matsuri hoảng hốt ôm đầu. Subaru nấp sau chăn. Pekora dừng giữa chừng chuẩn bị đòn tiếp theo. Fubuki suýt quăng cả ``bom gối'' vào Mio.

``K-Kuon-senpai?'' Fubuki ngoái nhìn người vừa thốt ra tiếng hét đó.

Trong ánh sáng mờ ảo từ Mặt trăn, khuôn mặt của cậu Tổng hắt lên. Giờ đây, đó không còn là một Kuon đang mệt mỏi muốn đi ngủ nữa, mà đó chính là gương mặt của một người đàn ông đã vượt quá giới hạn chịu đựng. Một loại sát khí kỳ lạ, vô hình mà ai cũng cảm nhận được, tỏa ra quanh người cậu như Stand trong J*J*.

Rushia nhìn khuôn mặt tối đen đó, vội lùi đằng sau mấy bước. ``T-Thôi xong rồi! Ta đã đánh thức anh ấy-nanodesu!''

``Anh ấy mà đã thốt lên bằng tiếng mẹ đẻ là đã thấy điềm rồi đó\dots'' Choco run rẩy. ``Lại còn cả sát khí kia nữa chứ\dots''

Nàng Cổ động nuốt nước bọt: ``T-Ta thực sự làm anh ấy điên lên rồi\dots''

Không đợi họ kịp chuẩn bị, Kuon lặng lẽ cúi người, nhặt lấy một chiếc gối. Đúng chiếc gối nặng trịch vừa ném qua cậu ban nãy.

Cả phòng đông cứng, mắt không rời khỏi tay cậu.

``C-Cái gối đó của em mà!'' Pekora ngạc nhiên. ``Cái gối đó chắc là không-''

Chưa kịp để cô nói hết câu, cậu Tổng đã dồn hết sức ném chiếc gối nhẹ hều kia.

*VÚT!*

Chiếc gối bay như tên bắn, phóng qua mặt của Towa và Nene khiến cho họ giật mình, ngã đáp mông xuống đệm.

``C-Cái gối kìa! Watame-chan!'' Mio hô hoán cảnh báo, nhưng đối phương còn chưa kịp phản ứng gì thì\dots

*BỤP!* ``Hự!''

Chiếc gối nhẹ như bông của nàng Thỏ trúng mục tiêu vào nàng Cừu đang ngồi co ro trong góc, giờ đây thì đã đổ gục xuống như bị đo ván.

Không chỉ thế, cú ném cực mạnh đó làm văng luôn cả cặp sừng cừu của cô trước những ánh mắt đầy ngỡ ngàng của tất cả mọi người.

``Cú vừa rồi mạnh quá p-p-peko\dots'' nàng Thỏ điếng người, nhìn thấy đàn em tội nghiệp bị hạ gục đầy tức tưởi.

``H-Hích\dots\ K-Kuon-senpai? Em nghĩ là a-anh có vẻ quá hăng trong việc này rồi đấy\dots'' Nene vừa cười vừa run, cố gắng hòa hoãn mọi thứ.

Chẳng thể ngờ, đó lại là câu nói cuối cùng của mình trước khi chiếc gối thứ hai phóng tới như đạn xoáy\dots

*BỤP!* ``Ặc\dots!''

Cú ném mạnh tới độ nàng ``Hải cẩu'' lăn ba vòng trên tấm nệm rồi ngã chổng vó lên trời, bất tỉnh nhân sự tại chỗ.

``N-Nene!'' Polka lao tới cô bạn với đôi mắt trợn ngược kia, cố lay dậy, trong khi Suisei toát mồ hôi hột: ``N-Này, Kuon-senpai à, em nghĩ anh nên dừng rồi đó\dots\ Bọn em sẽ không làm phiền anh nữa đâu!''

Nhưng cậu Tổng vẫn để ngoài tai ba người họ. Cậu cúi xuống, nhặt thêm hai chiếc gối nữa, rồi nói với giọng đáng sợ: ``Mấy đứa muốn chơi ném gối hả?''

Cậu hạ giọng xuống lần nữa, gần như gằn từng chữ một: ``Vậy anh sẽ cho mấy đứa biết thế nào là bị bóng đè!''

Thế là, từ trận ném gối đầy vui nhộn đột ngột trở thành một trận chiến sinh tồn.

Ngay sau câu nói đáng sợ của kẻ vừa mới ngủ dậy đó, các cô gái đồng loạt chạy tán loạn, tìm cách thoát khỏi sự truy lùng của Kuon. Người chui dưới bàn, kẻ núp sau nệm, có người tìm luôn A-chan để nấp phía sau, chỉ để nhận ra cô nàng vẫn đang ngủ ngon lành như trốn nợ.

``Đ-Đúng là khi anh ấy tức điên thì quả thật là đáng sợ\dots'' Miko run rẩy, áp sát vào người Suisei. Hai người họ đang núp dưới đống chăn màn.

``Nhất là khi làm phiền đến giấc ngủ của anh ấy nữa\dots'' cô bạn đáp lại. ``Pha vừa rồi là hơi dại đó\dots''

``K-Kuon-senpai! Bọn em không cố ý thật mà!'' Polka ngẩng đầu lên khỏi nơi ẩn nấp, trước khi bị Botan dìm xuống, đúng lúc một chiếc gối từ tay của cậu Tổng ném thẳng về phía cô.

``\hl{Đồ ngốc! Bà muốn bị anh ấy giết chết hả!?}'' nàng Sư tử trắng gắt lên khẽ, đồng thời suỵt một tiếng.

``\hl{N-Nhưng tôi chưa bao giờ thấy anh ấy lại hung bạo như thế cả!}'' nàng Cáo sa mạc lắp bắp, mắt ngước về cậu con trai đang đi dò khắp phòng như một con thây ma.

Tất cả mọi người dều đang nấp ở sau những màn chăn, màn gối, không dám ho he đến một tiếng.

Ngoại trừ đúng một người đen đủi vẫn còn đang nằm ngủ say như chết, không hề nề hà gì đến mọi thứ xung quanh.

``S-Sora-senpai! Chị ấy vẫn ở ngoài kia!'' Polka trố mắt nhìn cô nàng Thần tượng vẫn đang say giấc.

``Đừng có dại thế chứ, Omaru!'' Lamy ghìm chặt cô bạn xuống chặt hơn.

Bầu không khí chợt trở nên ngột ngạt hơn. Và rồi trong một diễn biến bất ngờ\dots

*VÚT!*

\dots một chiếc gối khác cũng được ném ra, mạnh như một viên đạn đại bác.

Lần này, mọi con mắt đều dõi theo đường bay của nó, và cả căn phòng cùng bật tiếng hét kinh hoàng ngay khi họ biết cậu đang nhắm tới đâu: ``\textbf{K-Khôôôôôông!!!!}''

Chiếc gối bay thẳng đến Sora, đúng, là Tokino Sora, người từng được coi là ``miễn nhiễm với mọi hiện tượng vật lý'' của \hll. Cô chỉ vừa ngồi dậy, mắt còn ngơ ngác sau tiếng ồn, chưa kịp cất tiếng chào.

``Ủa\dots? Có chuyện gì mà mấy đứa-''

*BỘP!*

Một cú đập hoàn hảo vào trán.

Nàng Thần tượng im bặt sau cú đáp của chiếc gối định mệnh đó. Một chiếc gối êm ái, không quá nặng, nhưng cũng đủ khiến cô sững lại. Mắt cô mở lớn trong thoáng chốc vì bất ngờ\dots

``A\dots''

\dots rồi chầm chậm khép lại như một thiên thần đang ngủ quên. Cô nàng ngã ra phía sau với vẻ duyên dáng, khiến cho bất cứ ai chứng kiến cảnh này cũng phải há mồm choáng ngợp.

Sora hạ cánh một cách nhẹ nhàng lên đống gối được dựng trước đó, tiếp tục chìm vào trong giấc mộng khuya.

Không ai hiểu điều gì vừa xảy ra.

``Đ-Đó là Sora-senpai\dots\ đúng không?'' Towa mấp máy, giọng vẫn còn run run.

``Kuon-senpai\dots\ vừa làm điều không tưởng kìa\dots'' Mel khựng lại.

``Chúng ta\dots\ có khi nào sắp bị xóa ký ức không đây\dots?'' Noel bồn chồn sau cú đập gối đó.

Không khí lạnh đi vài độ, khiến căn phòng điều hòa mát rượi dần hóa trở thành một cái tủ đông.

Kuon vẫn đứng đó, tay hạ xuống, đôi mắt như ánh lên một tia sét ngắn. Nhưng rồi, như có điều gì đó chạm vào bản tâm của mình, cậu hít sâu, bình tĩnh trở lại, nhìn về phía Sora đã ``bị hạ đo ván''.

``\dots Xin lỗi em, Sora-chan,'' cậu nói nhẹ như gió thoảng, thở dài.

Bất chợt, họ cảm thấy có một cảm giác nào đó lành lạnh chạy dọc sống lưng khi cậu từ một con gấu săn mồi bỗng chốc hóa về một người bình thường một cách nhanh đến chóng mặt.

``A-Anh ấy vẫn là anh ấy đấy chứ?'' Nene ngó nhìn cậu Tổng đang nhặt lại chiếc gối, nhưng dường không có ý định ném đi nữa.

``H-hình như là vậy\dots'' Lamy cũng run rẩy.

Cậu quay về con người bình thường của mình, nói với một giọng uể oải vì quá mệt: ``Thôi, đủ rồi. Một giờ mười lăm sáng rồi đấy. Mọi người, đi ngủ thôi.''

Một giọng nói nhẹ nhàng như mọi khi, nhưng lại đầy sức nặng tới mức không ai phản đối. Thậm chí những kẻ nghịch ngợm nhất trong những người còn thức kia còn không dám cựa quậy trái ý, vì chẳng ai biết được Kuon sẽ lại quay về bản tính hung dữ ấy.

Cả căn phòng đồng thanh khe khẽ: ``V-Vâng\dots!!''

Từng người từng người rón rén bò về nệm, kéo chăn kín đầu. Không ai hé lời, không ai dám hỏi thêm bất kỳ điều gì.

Có chăng chỉ là tiếng thở phì phò của người nào đó vẫn còn bị ám ảnh bởi khoảnh khắc ``Sora-senpai gục ngã''.

A-chan thì vẫn ngủ. Ngon lành như chưa từng có một trận chiến huyền thoại nào vừa xảy ra.

Cả căn phòng dần chìm vào trong sự yên ắng, không còn một tiếng xì xào nào nữa.

\pov{Hikawa Kuon}

\dots Tôi vẫn không ngủ được.

Đèn đã tắt, chỉ còn lại ánh sáng thoáng qua từ mặt trăng rọi vào, nơi mà tôi có thể nhìn rõ từ gương mặt của mọi người.

Tôi mở điện thoại lên xem đồng hồ. Hơn một rưỡi sáng. Chín mươi phút kể từ khi ngày sinh nhật của tôi đã chính thức hạ màn.

Tiếng xào xạc khe khẽ của chăn gối được kéo lên kín đầu. Tôi nằm dài giữa đống nệm, đầu vẫn còn hơi ong ong sau cơn thịnh nộ vừa rồi.

\dots Thật ra tôi cũng không nghĩ mọi chuyện lại đi xa đến thế. Tôi chỉ muốn ngủ. Chỉ là muốn được nằm yên, nhắm mắt, nghỉ ngơi sau một ngày dài trong niềm hân hoan của bữa tiệc sinh nhật. Nhưng rồi, hết gối này đến gối khác bay vào mặt\dots\ đến cả Sora-senpai cũng bị vạ lây bằng một cú ném thiếu kiềm chế từ tôi.

May mà không ai bị thương. Những cú đập gối cực mạnh của tôi vào các cô gái cũng chỉ dẫn họ tới những cơn say giấc nồng - mà thứ thực, có khi họ lại là những người may mắn hơn cả.

Tôi liếc nhìn quanh căn phòng đó một lần nữa.

Các cô gái nhà ta đều đã rúc vào chăn, ngoan ngoãn và im thin thít như những học sinh sau giờ phạt. Cái không khí căng như dây đàn vừa nãy giờ đã dịu xuống, nhưng tôi vẫn cảm tưởng rằng còn đâu đó vài ánh mắt len lén nhìn tôi như kiểu: ``Liệu anh ấy còn giận không nhỉ\dots''

Tôi chỉ có thể thở dài, rồi bật cười trước tình huống có phần quá lố vừa rồi.

Tôi chẳng hề giận ai cả. Mà đúng hơn, tôi thấy hơi buồn cười, vì bản thân đã có một pha có thể gọi là dọa nạt các thành viên.

Khi bầu không khí kia tan dần, vài tiếng cười rúc rích bắt đầu vang lên từ góc phòng, dẫu cho chúng có nhỏ, ngắn, như sợ ai đó sẽ lại nổi giận. Nhưng tôi chẳng nói gì, chỉ nhắm mắt lại, để mọi người tự điều chỉnh nhịp thở, tự tìm lấy khoảng lặng cho riêng mình.

\dots

\dots

Tôi vẫn chưa thực sự chìm vào giấc ngủ.

Tôi ngẩng đầu dậy, hé mắt một lần nữa.

Peko-chan đang nằm cạnh Subaru-chan, cả hai đã ngủ. Fubu-chan ôm gối quay mặt vào tường, đuôi tai cử động theo từng nhịp thở. Watame-chan vẫn còn nằm im sau cú ``hạ đo ván'' hồi nãy\dots\ Tôi hy vọng em ấy sẽ chỉ ngủ say thôi. Tôi không muốn vào nhà đá bóc lịch đâu.

Và đằng kia là Sora-chan. Mặt em ấy bình thản như thể chưa từng có gì xảy ra. Tôi mím môi, hơi chột dạ vì cú ném trời giáng đó. Mà, hình như em ấy chưa kịp tỉnh thì bị ăn trọn quả gối đó nhỉ? Ước rằng em ấy sẽ không hỏi tôi gì về chuyện này\dots

Dù vậy, sau những sự kiện điên rồ đó, một cảm giác ấm áp lan tỏa dần trong ngực tôi. Cả một căn phòng đầy những thành viên nhà \hll\ mà tôi đã làm chung trong vòng hơn hai năm, cả một đại gia đình ấm cúng, dù chỉ là tạm thời. Dẫu cho họ có hỗn loạn, dù có cãi cọ, thậm chí còn bị ``ăn gối'' thì họ vẫn ở đây, vì tôi. Không phải với tư cách là nhân viên hay cấp trên gì, mà là một người bạn.

Đúng, là chính tôi. Hikawa Kuon này đây.

Tôi bật cười khẽ, ngồi dậy lần cuối trong đêm nay.

Trước mắt tôi là hơn ba mươi con người đang say giấc nồng, những giấc mơ bình yên, những tiếng thở đều đều vang lên như bản giao hưởng của tình thân trong đêm giữa mùa hạ, thi thoảng còn văng vẳng những tiếng ve kêu ngoài kia. Có người rúc vào nhau, có người ngủ co như con mèo nhỏ, có người vẫn còn lẩm bẩm vài câu nói mớ.

Tôi thì thầm, như sợ phá vỡ sự yên lặng ấy: ``\hl{\dots Cảm ơn vì bữa tiệc sinh nhật độc nhất vô nhị này nhé\dots{} mọi người à.}''

Tôi nằm xuống trở lại một lần cuối, kéo chăn lên ngang ngực, để mặc bản thân trôi dần vào bóng tối và yên tĩnh.

Để rồi, bữa tiệc sinh nhật đặc biệt nhất trong cuộc đời tôi, bữa tiệc sinh nhật độc nhất vô nhị mà cả đời tôi có thể sẽ không thể trải nghiệm được lần thứ hai\dots\ thật sự đã khép lại, bằng một nụ cười và một giấc mơ ngọt ngào.

``\textit{Giá như bầu không khí này có thể kéo dài mãi mãi nhỉ\dots}''

Nhưng đó cũng chỉ là một giấc mơ hão huyền và có phần ích kỷ. Giờ tôi, trước mắt tôi sẽ còn là những ngày tháng gắn bó với công ty ``hề đội lốt thần tượng'' này dai dẳng đây\dots

\dots

\dots

\begin{center}
  \uline{Sáng hôm sau\\5:00 sáng.}
\end{center}

Ánh nắng đầu ngày chưa kịp rọi qua khung cửa kính mờ sương mùa hạ, nhưng đồng hồ báo thức của riêng tôi đã vang lên bằng một tiếng rung nhẹ trong túi áo.

``\vie{Sáng rồi cơ à\dots\ Sớm vãi\dots}''

Tôi ngồi dậy, vội tắt chiếc báo thức đi, rồi ưỡn vai lên cao, ngáp một hơi thật dài. Tôi nhìn quanh căn phòng, giờ đây vẫn là quang cảnh như đêm vừa rồi, nhưng có đôi phần lộn xộn hơn.

Cả gian phòng vẫn ngập trong giấc ngủ. Những dáng người cuộn tròn trong chăn, vài mái tóc rối xù lên trên gối, có người vẫn lẩm bẩm nói gì đó trong mơ, số khác ôm gối ôm như báu vật. Không ai hay biết trời đã dần sáng, và một ngày Chủ Nhật mới đang bắt đầu.

Phải thật rón rén\dots\ Nhẹ nhàng, lướt qua mọi người\dots\ Không để ai biết tôi đã dậy\dots

Tôi bước nhẹ, từng bước một, luồn qua giữa những chiếc chăn và đống gối còn vương vãi từ đêm trước. Mỗi bước đi là một thử thách giữ thăng bằng và tránh né, tới mức có một lần tôi suýt giẫm lên cái đuôi cáo mềm mềm của Fubuki, người mà tôi đang có cảm giác vẫn còn đang dỗi sau trò đua xe đó.

Tôi ra khỏi phòng mà không đánh thức bất kỳ ai, chạy khẽ khàng lên tầng năm. Căn nhà vẫn chìm trong tĩnh mịch, chỉ có tiếng bước chân tôi hòa cùng tiếng chim líu lo từ bên ngoài căn nhà.

Mọi người đều cần được ngủ thêm, tôi nghĩ thế. Dù gì tôi biết họ đã bỏ ra rất nhiều công sức để tổ chức cho tôi một buổi tiệc sinh nhật đầy hoành tráng này, nên họ thực sự xứng đáng có một giấc nghỉ dài.

Tôi rửa mặt, pha một cốc cà phê, rồi bắt đầu soạn vài ghi chú trong điện thoại, bản nháp cho một lịch trình nho nhỏ của ngày Chủ Nhật.

Vì, rốt cuộc thì\dots\ các bạn biết rồi đấy. Hôm nay vẫn là Chủ Nhật.

Và Chủ Nhật thì nghĩa là gì?

Một sự kiện điên rồ nào đó sẽ lại tiếp tái xuất ngày tại văn phòng của \hll\ chúng ta, lần này có sự góp mặt trực tiếp của ba ``tên tuổi'' đầy hứa hẹn: Akutan, Kana-chan, và Koro-chan.

Tôi ngẩng đầu lên, nhìn ra ngoài trời, khẽ thở dài chán chường, lẩm bẩm: ``\eng{Chết tiệt. Lại chuẩn bị lần nữa rồi.}''

Một tuần mới sẽ lại bắt đầu. Một kỷ nguyên tiếp theo của sự điên rồ, hỗn loạn, và tiếng cười đã sẵn sàng mở màn.

Nhưng đó là chuyện tôi sẽ kể ở *Quyển số Chín: Giữ vững chí khí của Cơn gió lốc mùa hạ!* Còn bây giờ, tôi đường như vẫn còn muốn tận hưởng khoảng thời gian ít ỏi cuối cùng trước khi thực sự bắt đầu vào thực chiến.

\ct

Buổi tối ngày hôm đó, sau ngày Chủ Nhật mà tôi phải giải quyết chuyện Baqua bị ``hồn lìa khỏi xác.'' Hôm nay đúng là một ngày kỳ lạ.

Tôi tiến về căn phòng tầng ba, nơi đã tổ chức một bữa tiệc cực lớn ngày hôm qua.

Tôi mở cửa và bước vào bên trong, bật chiếc đèn lên. Quả đúng như Akutan và Koro-chan nói, nơi này đã được dọn dẹp sạch sẽ, không còn nhìn thấy sự hiện diện của một buổi tiệc nữa, mà nó đã trở thành một căn phòng trống bình thường\dots

\dots tất nhiên, ngoại trừ tác phẩm mà chính tay ba mươi hai cô gái cùng gia đình tôi đã khắc họa nên. Vẫn là những nét vẽ đó, những mảnh giấy kia, những đồ vật trang trí đó, chúng dường như được tinh chỉnh lại một chút sao cho ngăn nắp hơn.

Bức tranh về tôi, nằm ở ngay chính giữa bức tường lớn nhất, còn được phủ rèm để trông như một tác phẩm nghệ thuật ở viện bảo tàng vậy.

Hây chà. Chưa gì tôi đã thấy nhớ bữa tiệc sinh nhật này rồi. Nhưng còn điều gì để hối tiếc đâu, đúng không? Tác phẩm do chính mồ hôi công sức và bao tấm lòng được gửi gắm vào từng sơn vẽ, từng giấy dán, nó vẫn sẽ ở đó, và có thể là trường tồn\dots

\dots nếu chúng tôi có thời gian tu sửa lại chúng thường xuyên. Thôi thì\dots

Tôi tắt đèn căn phòng tầng ba, nơi tụ họp của các thành viên \hll\ lại, cúi đầu xuống một lần cuối như là một lời tri ân, rồi rảo bước lên tầng năm, nơi mà gia đình tôi đang chờ đợi.

``\vie{Anh về rồi đây\dots}'' tôi cất tiếng gọi.

Đập vào mắt tôi là cô em gái đang ăn miếng bánh kem sô-cô-la ngon lành, một nụ cười mãn nguyện nở trên môi. Em ấy ngồi xếp bằng trên sàn, cười toe với chiếc đĩa rỗng trước mặt như thể vừa giành được chiến tích.

Trong khi đó, mẹ tôi thì vẫn đang chuẩn bị bếp trên kia, còn bố tôi thì vẫn chưa về. Dạo gần đây ông ấy cũng hay bận việc mà, cũng chẳng trách được.

Tôi thì có lẽ đã bắt đầu ngán bánh kem rồi, nên tôi cũng không tỏ vẻ ganh đua khi cô em gái ngốc nghếch đánh chén nốt miếng cuối cùng trong sự hạnh phúc.

``\vie{A, Onii-san! Mừng anh về\~{} Em lỡ ăn hết bánh kem rồi\~{}}''

``\vie{Thì anh có nói gì đâu,}'' tôi đáp lại, tự rót cho mình một tách trà ấm, dự định sẽ kết thúc ngày hôm nay một cách bình lặng, như bao ngày khác.

``\vie{Mà em ăn thế không phần lại cho Coco-chan với Aloe-chan à?}''

``\vie{Có chứ! Em vẫn còn để lại hai miếng cho các chị ấy đó!}''

``\vie{Thế thì tốt rồi. Tưởng họ sẽ chết đói chứ,}'' tôi đùa.

Ngay khi vừa dứt câu, hai giọng nói của hai người con gái vang lên từ hành lang cầu thang. ``Bọn chị về rồi đây\~{}!''

Cánh cửa ra vào bật mở. Vừa mới nói xong, nàng Rồng và nàng Succubus lại xuất hiện, với\dots\ ờm, cái gì đây?

``Bọn em đi vắng mấy ngày, Kuon-paisen thấy nhớ chứ?'' nàng Rồng cam gầm lên, tỏ vẻ oai phong với bộ dạng có phần kỳ quặc này của cả hai người.

``Về rồi à?'' tôi cố gắng giữ bình tĩnh, đặt cốc trà xuống. ``Chuyến đi của hai đứa suôn sẻ chứ?''

``Tuyệt, tuyệt lắm chứ! Ở bên xứ cờ hoa có thứ gì cũng to gấp đôi cả!'' nàng Rồng cười khoái chí, nhấc bổng chiếc vali nhẹ như bẫng, trên lưng còn đeo một chiếc túi vải to bảng, bên ngoài dán đầy sticker hình rồng uốn lượn mặc đồ cao bồi.

``Dám chắc là em đi hết mọi nơi ở xứ cờ hoa đó sao?'' tôi nhíu mày, không mong đợi gì từ câu trả lời của em ấy, vì chỉ cần nhìn những gì Coco-chan phô trương là tôi cũng đoán được phần nào.

Ở phía bên kia, Aloe-chan cũng kỳ quái không kém: em ấy ôm chặt một các-tông cũ, trên đó được dán chữ ``DO NOT OPEN AT NIGHT'' (không được mở vào ban đêm) bằng nét bút đỏ to tướng (tôi đoán em ấy định nặn ra một câu gì đó thật ngầu lòi bằng tiếng Anh, nhưng em ấy - cũng như đại đa phần các thành viên khác - hoàn toàn mù tịt ngoại ngữ nên mới đẻ ra cái câu nghe không ăn nhập kia). Trên vai của em còn vắt một tấm khăn lông hồng chóe với chữ ``SEXY DEMON'' thêu kim tuyến, mà có vẻ là quà lưu niệm từ buổi diễn nào đó.

Aloe-chan liếc quanh, nhướng mày: ``Ủa, không ai ra đón tụi này luôn hả? Tệ ghê.''

``Cái kiểu ăn mặc kỳ quặc kia làm bọn anh tưởng là trộm không đấy\dots'' tôi đáp. ``Mà mấy đứa vừa đi đâu về vậy?''

``Ủa, anh không biết sao, senpai?'' nàng Rồng đặt những đồ lỉnh kỉnh xuống, rồi tiến tới ngồi xuống sàn, ``Em vừa tới bên kia để thăm họ hàng đấy!''

``Còn ta vừa có một chuyến lưu diễn nhỏ quanh thành phố này!'' nàng Succubus cũng theo đó tiếp lời, đặt chiếc hộp lọc xọc đồ đạc xuống bên cạnh. ``Ryuuhou? Không chào ta một tiếng à?''

``Bọn em đang ăn bánh còn sót lại từ tiệc sinh nhật của Onii-san đó!'' cô em của tôi reo lên, tay giơ lên miếng kem sô-cô-la cuối cùng trên dĩa lên và thưởng thức đến cùng.

``Tiệc sinh nhật?'' cả hai người họ đồng thanh, rồi cùng nhìn thẳng sang tôi.

``Vâng!'' Ryuuhou nhanh nhảu kể, nét mặt hớn hở như được giao nhiệm vụ tối mật. ``Hôm qua bọn em tổ chức tiệc sinh nhật cho anh ấy! Đông vui lắm luôn! Bạn của Onii-san kéo tới đầy nhà luôn, còn ngủ lại nữa đó!''

``Bạn của Kuon-paisen\dots?'' Coco ngước lên, hơi nghiêng đầu. ``Ý em là mấy người quen ở đại học vói hồi trung học ấy hả?''

Tôi khẽ gật đầu, vẫn giữ vẻ bình thản và từ tốn giải thích: ``Ừ. Mà thề, mấy ông ấy tự lên kế hoạch hết, anh chẳng biết gì cho tới khi bị bố anh mời xuống tầng luôn.''

``Nó như thế nào vậy?'' Coco tiếp tục.

``Nói ra khá là dông dài, nhưng mà bọn anh cũng có khá nhiều trò vui ở đó,'' tôi nhún vai.

``Nói toẹt thẳng ra đi cho em nghe nào,'' Aloe tỏ vẻ bắt đầu mất kiên nhẫn, tỏ vẻ muốn được nghe câu chuyện nửa thật nửa thêu dệt cho hai anh em tôi ``bất đắc dĩ làm đạo diễn''.

``Ôi, nhiều lắm!'' cô em tôi bật ngón tay cái. ``Bất đầu từ hát karaoke, rồi đua xe, thậm chí còn lác hẳn sang đập dưa nữa!''

``Nghe có vẻ thú vị đấy! Kuon-paisen chắc là cũng khá vui vẻ với đám con trai bọn anh đấy nhỉ?''

``Cũng\dots\ có thể hiểu như vậy. Và rồi ngay sau đó họ rủ anh chơi thử thách hay sự thật, trò kể chuyện ma, và kết thúc bằng trò ném gối.''

Mẹ tôi từ trên bếp xuống, cũng gật đầu phụ họa: ``Đúng đấy. Cô còn nhớ lúc cô chú test mic, có con bé tóc hồng- à nhầm, bạn của thằng Cò bật nhạc sớm, thế là cuối cùng tất cả mọi người đều cười lăn lộn, còn cậu ta thì đỏ mặt tới mức hờn dỗi luôn!''

Tôi ho khẽ, trừng mắt nhẹ về phía bà ấy. Đây không phải là lần đầu mẹ tôi vô tình thốt ra những chi tiết gây bất lợi như thế. Tôi vẫn còn nhớ như in hồi mẹ tôi nói oang oang với các bác sĩ ở bệnh viện gần nhà rằng tôi bị rò hậu môn trong khi đang chữa đường ruột - lúc đó là tầm tôi chuẩn bị thi tốt nghiệp trung học - để rồi cuối cùng phải chuyển viện gấp trong ngày chỉ để ``trốn'' phẫu thuật đấy.

``Nói tóm lại thì, ngày hôm đó y như là một lễ hội á!'' em gái tôi cười giả lả, tìm cách đánh lạc hướng họ khỏi chi tiết ``lỡ miệng'' kia. ``Coco-onee-san với Aloe-onee-san hôm đó vắng nhà tiếc thật đấy!''

Nàng Succubus nheo mắt nhìn tôi, lẩm bẩm: ``Hừm\dots\ Mà sao em thấy hình như có gì mờ ám ở đây nhỉ? Anh giấu gì sao, Kuon-kun?''

Tôi quay đi, nhấp một ngụm trà. ``Thực ra\dots\ hai đứa không có ở nhà lại là điều may đấy. Chứ nếu có mặt, anh cũng chẳng biết phải ăn nói thế nào với Tanigo-san nữa\dots''

Hai cô nàng ở đậu liếc nhau. Không cần nói ra, cả hai đều hiểu rõ nếu tôi nhắc đến tên ``\emph{Tanigo-san}'' - người đại diện cao nhất - thì chắc chắn cái bữa tiệc kia không đơn giản chỉ là vài người bạn tụ tập ăn bánh sinh nhật.

Thấy rằng họ vừa động phải một chủ đề khó nói, nàng Rồng gãi má, rồi lái hẳn sang chủ đề khác: ``À, nhắc mới nhớ! Em có mua quà lưu niệm cho mọi người đây!''

``Em cũng có đồ mang về chia mọi người nữa này!'' nàng Succubus tiếp lời, mở chiếc hộp bìa các-tông ra.

``Chẳng trách mà hai đứa mang về đống đồ lỉnh kỉnh thế,'' tôi chống hông, bật cười nhẹ. ``Được rồi, hai đứa mang cái gì về, đưa đây cho anh xem nào.''

Coco-chan lôi ra trong túi xách đeo vi của mình mấy món đồ trông\dots\ chẳng ăn khớp gì nhau lắm: một chiếc bật lửa hình đầu lâu biết nháy mắt, một lọ nước hoa có mùi\dots\ thịt nướng BBQ, và một bức tượng nhỏ hình rồng mặc bikini.

``\dots Ở xứ đó có bán những vật thể kỳ quặc như thế sao?'' tôi nhướng mày, không biết phải phản ứng như thế nào.

``Đúng rồi!'' em ấy đáp, đặt ba món đồ đó lên bà làm việc của tôi. ``Tặng anh cái này, để bàn ở đây là hợp luôn.''

``Ờm\dots\ Cảm ơn?'' tôi cất giọng đầy khó hiểu. ``A-Anh cũng không biết nên nhận tấm lòng thành này của em ra sao nữa\dots''

Aloe-chan cũng đặt cái hộp bìa lên, cắt phần băng dính ra, để lộ một đống huy hiệu phát sáng đủ màu, một chiếc băng đô có chữ ``I AM A STAR'' (Ta là ngôi sao) lấp lánh, và một chiếc loa mini hình\dots\ quả chuối.  Điện thoại hình quả chuối tôi có thấy qua rồi, nhưng chiếc loa hình này lại là\dots\ một thứ gì đó khang khác.

``Chọn đi, quà này là phần thưởng từ fan đó,'' em ấy hớn hở nói. ``Một mình ta có lẽ không dùng hết chỗ này đâu, nên là cứ lấy thoải mái!''

``Vậy em lấy cái băng đô'' em tôi nhìn vào đống quà, hí hửng lấy món quà đó. ``À, với cả mấy cái huy hiệu cài áo nữa!''

``Anh thì\dots\ chắc là lấy luôn cái loa nhỉ?'' tôi vớ lấy chiếc loa mini. ``Hình thù hơi dị, nhưng ít nhất cũng tiện dụng phết.''

Đúng lúc này, mẹ tôi vừa từ bếp xuống lần nữa, trên tay là hai miếng bánh kem - một vị dâu, một vị cà phê - còn sót lại từ chiếc bánh khổng lồ kia. ``Các cháu có muốn ăn không? Bạn của thằng Cò làm đấy.''

``Vâng\~{}''

Cả hai cô nàng sau đó vui vẻ ngồi lại với chúng tôi kể lại những gì họ đã đi và họ đã làm trong suốt hai ngày đó, rồi cùng gia đình chúng tôi dùng bữa.

Nói gì thì nói, sự trở về của Coco-chan và Aloe-chan khiến căn nhà lại trở về nhịp thường ngày, một chút ồn ào, náo nhiệt, nhưng đầy hơi ấm. Dù có đôi lúc tôi ước được yên tĩnh một chút, thì chính cái sự hỗn độn này mới khiến nơi này trở thành nhà.

Thật lòng mà nói, cái ``ổ'' này\dots\ mới là nơi khiến tôi thấy mình thực sự thuộc về.

\ct

Sau khi dùng bữa tối xong, Coco-chan cầm một lon soda trông khá màu mè, vừa khui vừa huýt sáo, thong dong quay trở về phòng mình, còn Aloe-chan thì ôm chiếc hộp bìa của mình lên.

``Thôi, em về phòng đây\~{}'' nàng Succubus ngân nga đầy phấn khích, lướt về phía cầu thang tầng 2 như thể sắp bước lên sân khấu, ngân nga trong giai điệu ``Ochame Kinou''.

Tôi nhìn theo bóng cô nàng, rồi quay lại ăn nốt miếng bánh kem còn dang dở. Ryuuhou ngồi kế bên, chống cằm như đang ngẫm nghĩ điều gì đó.

``Em nghĩ\dots\ nếu hai người họ mà về sớm hơn một hôm, chắc sẽ được tham gia bữa tiệc sinh nhật của anh rồi nhỉ?'' cô bé lên tiếng.

Tôi thở dài một tiếng. ``\vie{Em quên mất là ta đang phải giữ bí mật chuyện Coco-chan và Aloe-chan ở nhà ta à? Yêu cầu của chú YAGOO-san đấy.}''

Cô em gái của tôi nghiêng đầu, định hỏi lại thì chợt nhớ lại giao kèo mà tôi nói với cả gia đình tôi nửa năm về trước. ``\vie{À, ừ, đúng rồi. Em quên.}''

``\vie{Thì đấy,}'' tôi đưa nốt miếng bánh cà phê mà Coco ăn không hết vào miệng, ``\vie{Họ mà có mặt thì chả biết cả mọi người ở công ty sẽ xử trí thế nào đâu.}''

Chúng tôi đang định tiếp tục trò chuyện linh tinh, bỗng một tiếng hét vang dội từ tầng 2 vang lên, chấn động cả trần nhà.

``\textbf{Cái quái gì thế này!? Chuyện gì xảy ra với căn phòng tôi đây!?}''

Tiếng la quen thuộc khiến Ryuuhou phụt nước, đứng bật dậy, mặt tái xanh, trông như thể vừa mới nhớ ra điều gì đó quan trọng. ``\vie{C-Chết rồi\dots!}''

Tôi nhìn em gái mình, mặt nghiêm lại: ``\vie{Sao đấy? Lại làm gì với phòng của em ấy à?}''

Em gái tôi mặt xanh lè, lắp bắp kể lại: ``\vie{C-Chả là\dots\ hôm qua Aqua-onee-san và Shion-onee-san với cả em mượn tạm phòng Aloe-onee-san để làm bánh kem sinh nhật anh ạ. Vì bếp tầng 6 thì bố dùng để chuẩn bị đồ nướng, còn tầng 4 mẹ lại đang nấu lẩu ạ\dots}''

``\vie{Và thế là mấy đứa kéo xuống dưới đó? Mọi người có biết đó là phòng Aloe-chan không?}''

``\vie{K-Không. Em nói dối rằng đó là phòng của em đấy ạ.}''

``\vie{Và sau khi làm bánh xong mấy đứa cũng không thèm dọn lại luôn?}'' tôi rướn mày.

``\vie{Thì thế, nhưng mà\dots\ bọn em cuốn vào bữa tiệc sinh nhật của anh quá, xong rồi quên mất luôn\dots}''

Tôi chỉ thở ra một tiếng, ngụ ý: ``Chán lắm cơ,'' lắc đầu rồi đứng dậy.

``Đi thôi. Xuống xem tình hình thế nào không lại để Aloe-chan lên nhà mình ngủ bây giờ.''

Hai anh em cùng chạy vội xuống tầng hai.

Cánh cửa phòng của nàng Succubus vẫn mở toang. Trước mắt chúng tôi là một cảnh tượng bừa bộn không lời nào diễn tả nổi: kem tươi trên tường khô lại thành mảng, bột mì phủ dày mặt sàn như tuyết, mấy cái giấy gói bột dùng xong nằm xoắn thành cuộn, mấy quả dâu bị bẹp dúm nằm rải rác như trận địa, một vết lõm rõ to trên trần nhà (mà theo lời của Ryuuhou là do Shion-chan hóa lớn chiếc lò), còn vữa tường thì bong từng mảng.

Nói tóm lại, trông không khác gì một chuồng lợn hay một căn phòng bị tàn phá bởi cơn bão.

Aloe-chan đang đứng giữa căn phòng với ánh mắt kinh hoàng, tay nắm chặt túi đồ lưu niệm như sắp bóp nát tới nơi.

Ryuuhou rụt cổ lại khi ánh mắt của em ấy lia sang, còn tôi thì vội giơ hai tay lên: ``A-Anh không biết gì hết nhá!''

Cô nàng Succubus gằn giọng, hỏi cô em gái tôi đang toát mồ hôi hột: ``Thế này là thế nào đây hả Ryuuhou? \textbf{Tại sao phòng chị lại trông như ổ gà ổ voi thế này!?}''

``E-Em xin lỗi! Hôm qua bọn em có mượn phòng để làm bánh, thế rồi bọn em quên mất phải dọn lại\dots\ Giờ em mới nhớ ra\dots''

Tôi đằng hắng một tiếng, chen vào trước khi cô nàng ương ngạnh đó nổi cáu thêm: ``Thôi, thôi được rồi, đừng cáu. Chúng ta sẽ cùng nhau dọn dẹp lại chỗ này, còn ý kiến ý cò gì thì tính sau. Tất nhiên là em cũng phải phụ bọn anh dọn nữa, càng nhiều người làm càng nhanh.''

Aloe khoanh tay, hất tóc một cái đầy bực bội: ``Em vốn dĩ không thích dọn phòng mình đâu đấy. Tại sao lại thế?''

Cô em gái của tôi thở dài: ``Thì đây là phòng của chị mà. Em thấy chị siêng dọn phòng của Onii-san với phòng bác Hijaki-san lắm?''

``Đó là phòng của họ, việc nào ra việc đấy! Với cả, nhìn phòng họ dính bụi, ta thấy ngứa mắt lắm!''

``Thế em không thấy phòng mình ngứa mắt à?'' tôi hỏi vặn lại. ``Phòng em giờ còn không hổ lốn bằng phòng anh, mà em vẫn dọn đấy thôi. Tính lấy cớ này để mon men sang phòng anh ngủ chung á? Không có đâu.''

Tới đây, Aloe-chan không biết lấy lời gì để bao biện vì bị tôi bắt bài. Cuối cùng, em ấy cũng chịu lùi bước, gắt lên: ``Đ-Được thôi! Dọn thì dọn! Nhưng mà bởi vì anh yêu cầu đấy nhé!''

Tôi chỉ mỉm cười nhẹ. Có vẻ cho dù là Aloe-chan ngang tàng và bề trên thật, nhưng em ấy vẫn không thể bỏ mặc mọi thứ trong tình trạng này, hoặc đúng hơn, không thể nào thao túng được tôi.

Cuối cùng, buổi tối hôm đó, tôi, Ryuuhou và Aloe-chan đã dành gần hai tiếng để dọn dẹp lại phòng ốc của cô nàng ở dậu này, giữa mùi kem bơ và vụn bánh quy vương vãi khắp nơi.

Cơ mà, ngẫm nghĩ lại, cũng lâu lắm rồi tôi mới động tay động chân dọn dẹp lại nhà cửa kể từ hồi Tết. Tôi không cuồng việc nhà như mẹ tôi, nhưng ít nhất cũng có việc gì đó để cho tôi làm ở nhà.

\begin{center}
  \uline{Sáng hôm sau\\Phòng 53-H1 - Trường Đại học Xây dựng Kawachi}
\end{center}

Tôi xách chiếc cặp máy tính lên vai, đeo tai nghe vào rồi nhét chiếc máy S*tch mới toanh vào ngăn phụ, mà tôi mới là lần thứ hai tôi động tay vào. Tôi vẫn chưa có thời gian để cài các trò khác hay kết nối mạng, nhưng chỉ riêng việc cầm nó đi học cũng khiến tôi cảm thấy\dots\ hơi sao sao. Tôi biết là tôi đã từng bắt gặp một sinh viên mang máy S*itch đi học ở thư viện năm ngoái, nhưng đến lúc tự tay cầm chiếc máy lên thì đó lại là một trải nghiệm khác.

Trên đường đến lớp, tôi vẫn còn ngáp ngắn ngáp dài. Thức dậy sau một đêm chạy đồ án môn học là không phải cách hay để mở đầu tuần mới. Nhưng ít nhất thì mấy trò hỗn loạn hậu sinh nhật cuối tuần đã kết thúc, mọi thứ đã quay trở về guồng quay sinh hoạt thường ngày.

Tôi vừa vào đến lớp thì đã thấy Yoisai vẫy tay từ xa: ``\vie{Ê, Kuon ơi! Biết ngay là ông sẽ tới lớp với cặp mắt gấu trúc mà.}''

``\vie{Đêm qua chạy KPI đồ án, buồn ngủ éo chịu được\dots}'' tôi cười nhạt đáp lại, kéo chiếc ghế ngồi xuống bàn ba người họ.

Cậu bạn Lập trình ngó nghiêng nhìn chiếc cặp của tôi, còn cậu bạn Đeo kính vừa xé gói bánh ăn sáng vừa thì thầm: ``\vie{\hl{Này, có phải ông đang mang cái máy S*itch không đấy?}}''

``\vie{Ừ. Máy S*itch chứ còn gì,}'' tôi gật đầu, lôi chiếc máy ra một chút như để xác nhận. ``\vie{Quà sinh nhật của tôi hôm vừa rồi.}''

Cô bạn Học bá của lớp cau mày: ``\vie{Nó là quà sinh nhật á? Ai tặng vậy?}''

``Sui-chan,'' tôi đáp gọn. ``\vie{Pol-pol với Senchou còn tặng nguyên cả thư viện game của con máy này nữa.}''

``Suisei, Polka với Marine à?'' cô bạn tròn mắt, ``\vie{Họ cũng chúc mừng sinh nhật ông sao?}''

Tôi tháo một bên tai nghe xuống, thở ra một tiếng rồi nhếch môi cười trừ: ``\vie{Không chỉ chúc miệng suông đâu. Họ còn làm hẳn cả một kế hoạch tổ chức sinh nhật bài bản cho tôi, mà đến phút cuối tôi mới biết đấy.}''

Ba người kia đồng loạt nhìn tôi như thể tôi vừa bảo rằng mình trúng xổ số giải đặc biệt.

``\vie{K-Khoan, cái gì!? Ông đang nghiêm túc đấy à!?}'' Yamazu lắp bắp.

Tôi gật đầu, chống cằm lên bàn, ánh mắt lơ đãng nhìn ra ngoài cửa sổ. ``\vie{Đê tôi kể nhá. Tối thứ sáu một ngày trước sinh nhật, A-san có gọi điện cho tôi hỏi thăm, rồi bảo tôi ngày hôm sau không phải đi làm.}''

``\vie{A-san bảo ông nghỉ một ngày á?}'' Yoisai nhíu mày. ``\vie{Sao lại thế?}''

``\vie{Hồi đó tôi còn chả biết là tại sao, mãi tới khi tối khi mà tôi xuống nhà thì thấy họ đã chuẩn bị sẵn một bữa tiệc sinh nhật dưới nhà, tôi mới luận ra là `à, hóa ra chị ấy bảo tôi nghỉ một ngày là vì thế'.}''

``\vie{Ra thế\dots}'' Tamachi đáp, vo lại gói bánh mới mở. ``\vie{Vậy trong bữa tiệc đó các ông có gì vậy?}''

``\vie{Nhiều trò lắm. Để tôi kể tiếp đây\dots}'' tôi ngồi lại với ba người, hồi tưởng lại những hoạt động vào buổi tối ngày thứ Bảy hôm đó. Họ đều chăm chú lắng nghe, từng chi tiết, từng bối cảnh, từng hoạt động được tôi kể rất kỹ, như thể chính họ cũng đang có mặt trong bữa tiệc hôm đó.

Yamazu há hốc mồm, rơi cả cái bình trong tay khi hiểu được mọi thứ: ``\vie{Vãi chưởng, thật luôn!?}''

``\vie{Đến tôi còn không ngờ tới luôn. Thề, ngay cả đến việc các cô gái \hll\ tổ chức sinh nhật tôi đã là điểu không tưởng rồi,}'' tôi ngẩng mặt lên, cười chát, ``\vie{đằng này lại còn làm tới mức này nữa chứ!}''

Tôi không biết mọi người đã chuẩn bị như thế nào, nhưng tôi chỉ biết răng họ đã làm rất tốt phần công việc được giao. Tôi tưởng tượng đến viễn cảnh Aqua-chan và Shion-chan đập nhau bằng âu bột trong phòng Aloe-chan, rồi khi bố mẹ tôi cùng mọi người nấu ăn ngon lành.

Hẳn là họ đã phải làm rất nhiều việc, và dù ban đầu tôi vẫn không hiểu tại sao họ lại làm như vậy, nhưng rồi tôi cũng đã hiểu ra rằng\dots

\dots tôi thực sự phải là một người cực kỳ, cực kỳ quan trọng tới mức nào để họ có thể bỏ chút nguyên một ngày để tổ chức một sự kiện trọng đại và hấp dẫn như vậy dành cho tôi.

Trong lúc tôi còn đang nghĩ về ngày sinh nhật hôm đó, bỗng chốc tôi nghe thấy tiếng bước chạy của các sinh viên trong lớp. ``\vie{Thầy! Thầy vào kìa!}''

Và thế là, cả lớp ai nấy đều quay về chỗ ngồi của mình, chuẩn bị cho một ngày mới. Bữa tiệc sinh nhật của tôi có thể đã trôi dần vào thời gian, nhưng những gì ấn tượng nhất vẫn sẽ được lưu giữ mãi trong lòng tôi, và cả tập thể các cô gái đã làm nên nó nữa.

\il

Tôi vẫn còn nhớ khá rõ buổi tối hôm thứ Sáu ấy, khi điện thoại của tôi reo lên và cái tên ``A-san'' hiện ra trên màn hình. Giọng chị dịu dàng nhưng cũng không giấu nổi sự nghiêm túc quen thuộc\dots

``Hikawa-san đúng không? Tôi có việc này muốn nói với cậu đây. Giám đốc chuyển lời cho cậu bảo ngày mai cậu không phải đi làm nữa nha.''

Tôi đã tưởng rằng chị ấy gọi để thông báo tôi bị đuổi việc, nhưng hóa ra chỉ là một lời khuyên bảo tôi ở nhà để tận hưởng ngày sinh nhật. Và dù tôi có một chút nghi ngờ, tôi vẫn răm rắp nghe theo.

Chỉ là\dots\ tôi không thể nghĩ đó lại là khởi đầu cho những gì theo sau.

Tôi đã nghĩ ngày hôm sau, ngày Mười bảy tháng Bảy, sẽ chỉ như bao lần sinh nhật trước, dù gì cũng đẫ là lần thứ hai mươi lăm rồi. Qua mỗi năm, niềm mong chờ lại vơi đi một chút. Hồi nhỏ thì háo hức không ngủ được, chỉ mong mau tới sáng. Lên cấp hai, cấp ba thì còn mừng vì có người nhớ. Còn những năm sau nữa khi đã bước chân vào cuộc sống đại học, tôi bắt đầu quen với một sinh nhật yên tĩnh: tự thưởng cho mình một chiếc bánh nhỏ, một tách cà phê, vài tiếng chơi game đơn độc và một bữa tối đơn giản với gia đình.

Năm đầu tiên tôi đi làm ở \hll, tôi thậm chí còn suýt quên mất ngày sinh nhật của mình, nhưng rồi bố mẹ cùng bác Hijaki-san đã tổ chức cho tôi một buổi tiệc nhỏ.

Giờ, tới mức mà cả các cô gái \hll\ cũng làm nguyên một bữa đại tiệc như thế\dots\ đó lại là điều mà tôi tưởng chỉ có trong anime hoặc là manga, giờ lại được thực hiện trong cuộc sống thật.

Họ đã lên kế hoạch với nhau từ trước. Tất cả đều giấu kín, không ai hé nửa lời. Và rồi từ lúc tôi bước qua cánh cửa căn nhà quen thuộc ấy, mọi thứ như bùng nổ.

Tôi bị ném vào một thế giới tràn ngập âm thanh và màu sắc: dây ruy băng bay loạn xạ, đèn nhấp nháy, một chiếc bánh khổng lồ (mặc dù một phần chiếc bánh hơi móp góc do tai nạn của Sui-chan), và hơn hết là tiếng cười nói rộn rã của các cô gái tôi luôn yêu quý, những ``người bạn'' tưởng như chỉ có thể kết nối qua cuộc sống công việc, nay lại đang ngồi vây quanh tôi ngoài giờ tan tầm.

Từ trò đập dưa trong phòng tầng ba, những cuộc đua xe hấp dẫn, bữa lẩu sôi sùng sục bên hiên, cho đến tiếng nhạc Miko lỡ tay bật nhầm giữa lúc thử âm thanh\dots\ và cả trận chiến ném gối đậm chất ``\textit{hologra}'' đến mức đèn tôi trở thành một người hạ gục được ``kẻ bất khả xâm phạm'' (theo lời của các cô gái) tất cả như một giấc mơ tuổi thơ mà tôi chưa từng có dịp được sống trọn vẹn.

Trong suốt bữa tiệc ấy, tôi nhiều lần tự hỏi: Mình thực sự xứng đáng với điều này sao?

Có thể là không.

Nhưng họ đã chọn ở bên tôi, cùng nhau tạo nên khoảnh khắc ấy, bằng cả trái tim.

Tôi từng nghĩ rằng sinh nhật là điều dành cho trẻ con, dành cho lứa tuổi nhi đồng thối tai. Nhưng giờ, tôi hiểu: không phải ai cũng may mắn được sinh ra, càng không phải ai cũng may mắn có người nhớ đến ngày mình được sinh ra.

Tôi đã sống hơn hai mươi tư năm trời mà không hề nghĩ rằng sẽ có một ngày, tôi được yêu thương đến thế, nếu không muôn nói là được tận tình đến tận răng.

Nếu tôi có một mong ước trong ngày sinh nhật mà tôi có thể cầu trong khoảnh khắc bài hát sinh nhật được ngân vang đó, thì\dots

\begin{center}
  ``Mong rằng, các thành viên \hll, cùng toàn thể mọi người\dots\ Từng người một chúng ta cũng sẽ có ngày được sống trong cảm giác ấy. Được tổ chức sinh nhật không chỉ bằng bánh kem và quà, mà bằng một vòng tay tập thể, nơi mọi người thật sự nghĩ cho nhau. Chúng ta hãy cùng vững bước và tiến về phía trước, để ngày sinh nhật này trở thành một kỷ niệm đặc biệt, một ngày đáng nhớ.''
\end{center}

Bởi vì, trong một thế giới mà ai cũng đang mải chạy theo lịch trình, chỉ cần dừng lại vì nhau một chút thôi\dots\ cũng đã là điều tuyệt vời lắm rồi.

Vậy đấy. Sinh nhật thứ hai mươi lăm.

Một ngày Mười Bảy tháng Bảy.

Một bữa tiệc mà tôi sẽ không bao giờ quên. Và có lẽ, kể cả khi tôi rời khỏi \hll\ một ngày nào đó trong tương lai\dots

\dots đó cũng sẽ là một điều gì đó mà tôi có thể đem ra để ôn lại kỷ niệm và nhớ lại.

Tạm biệt nhé, sinh nhật thứ hai mươi lăm của cuộc đời tôi.

\end{document}